\section{Header 1}
\lipsum[1]

\subsection{Header 2}
\lipsum[2]

\subsection{Header 3}
\lipsum[3]

\begin{theorem}{Gauss-Markov}{thm:gaussmarkov}
In a linear regression model $Y = X\beta + \varepsilon$ with assumptions:
\begin{itemize}
    \item $\mathbb{E}[\varepsilon] = 0$ (zero mean errors)
    \item $\text{Var}(\varepsilon) = \sigma^2 I$ (homoscedasticity and no autocorrelation)
    \item $X$ is full rank
\end{itemize}
the ordinary least squares (OLS) estimator $\hat{\beta} = (X^TX)^{-1}X^TY$ 
is the Best Linear Unbiased Estimator (BLUE), meaning it has the smallest 
variance among all linear unbiased estimators.
\end{theorem}

\begin{definition}{Gauss-Markov}{def:gaussmarkov}
In a linear regression model $Y = X\beta + \varepsilon$ with assumptions:
\begin{itemize}
    \item $\mathbb{E}[\varepsilon] = 0$ (zero mean errors)
    \item $\text{Var}(\varepsilon) = \sigma^2 I$ (homoscedasticity and no autocorrelation)
    \item $X$ is full rank
\end{itemize}
the ordinary least squares (OLS) estimator $\hat{\beta} = (X^TX)^{-1}X^TY$ 
is the Best Linear Unbiased Estimator (BLUE), meaning it has the smallest 
variance among all linear unbiased estimators.
\end{definition}

\begin{example}{Gauss-Markov}{ex:gaussmarkov}
In a linear regression model $Y = X\beta + \varepsilon$ with assumptions:
\begin{itemize}
    \item $\mathbb{E}[\varepsilon] = 0$ (zero mean errors)
    \item $\text{Var}(\varepsilon) = \sigma^2 I$ (homoscedasticity and no autocorrelation)
    \item $X$ is full rank
\end{itemize}
the ordinary least squares (OLS) estimator $\hat{\beta} = (X^TX)^{-1}X^TY$ 
is the Best Linear Unbiased Estimator (BLUE), meaning it has the smallest 
variance among all linear unbiased estimators.
\end{example}